% \noindent \textbf{Problem Formulation.} Given a complex reasoning task that requires multi-step reasoning and the retrieval of external knowledge to derive solutions, the goal is to produce, for each question $q$, a comprehensive solution consisting of a multi-turn reasoning chain $R$, the retrieved information $I$, and the final answer $A$. Specifically, for retrieval process, we consider two distinct retrieval infrastructures: dense RAG and graph-based RAG. 

% \noindent \textbf{Knowledge Graph Construction.}

% \noindent \textbf{Graph-based Retrieval-augmented generation.}

% \noindent \textbf{Retrieval-augmented generation.}

% \noindent \textbf{Answer Generation. }
We consider open-domain question answering with an input query $q$ and a large language model (LLM) $\mathcal{M}$. Under standard LLM reasoning, the model generates an answer $y$ conditioned solely on the query: $y \sim \mathcal{M}(q).$
While powerful, this formulation relies entirely on the knowledge encoded in model parameters and is limited in its ability to handle knowledge-intensive or multi-hop reasoning tasks that require access to external information.

\noindent\textbf{Retrieval-Augmented Reasoning (RAG).}
RAG extends LLM reasoning by grounding generation in an external knowledge corpus $\mathcal{K}=\{d_1,\dots,d_N\}$. A RAG system consists of a retriever $\mathcal{R}$ and an LLM $\mathcal{M}$. Given a query $q$, the retriever selects a set of relevant documents or text chunks $C_q = \mathcal{R}(q \mid \mathcal{K})$,
which are appended to the model input for answer generation: $y \sim \mathcal{M}(q, C_q).$
In most existing RAG pipelines, retrieval is performed once prior to decoding, and the retrieved context remains fixed throughout generation. This \emph{static} or \emph{one-shot retrieval} assumption underlies the majority of dense-retrieval-based RAG systems.

\noindent\textbf{Graph-Based Retrieval-Augmented Generation (GraphRAG)}
GraphRAG further extends retrieval-augmented reasoning by explicitly organizing the knowledge corpus into a structured graph or hypergraph prior to inference. We abstract a GraphRAG knowledge base as $\mathcal{G} = (\mathcal{V}, \mathcal{E})$,
where nodes $\mathcal{V}$ represent textual units (e.g., documents, entities, or summaries) and edges $\mathcal{E}$ encode relationships among them. Retrieval in GraphRAG typically selects nodes, paths, or subgraphs from $\mathcal{G}$, which are then provided as structured evidence to the LLM: $y \sim \mathcal{M}(q, Z_q), \quad Z_q \subseteq \mathcal{G}$. Compared to dense RAG, GraphRAG introduces explicit structural inductive bias that can improve multi-hop reasoning and evidence aggregation. However, similar to standard RAG, GraphRAG retrieval is typically in a single preprocessing step, and the selected evidence is appended to the model input before decoding.

In this work, we focus on benchmarking RAG and GraphRAG under \emph{agentic search} settings, where retrieval is integrated into inference and performed iteratively during decoding. This paradigm enables multi-round retrieval and adaptive information access based on intermediate reasoning states, fundamentally changing how retrieval interacts with LLM reasoning.
%Unlike the standard RAG or GraphRAG based LLM reasoning systems above, in this work, we focus on benchmarking RAG and GraphRAG under agentic search settings, where retrieval is performed iteratively during inference rather than appended as a one-shot input. This paradigm allows LLMs to perform multi-round retrieval and adapt information access based on intermediate reasoning states.
%Despite their differences in retrieval structure, both RAG and GraphRAG share a common assumption: retrieval is static and external to the decoding process. A fixed set of retrieved evidence is selected prior to generation and remains unchanged throughout inference. In contrast, recent agentic search systems integrate retrieval into the inference process itself, allowing LLMs to perform multi-round retrieval and adapt information access based on intermediate reasoning states. In this work, we focus on benchmarking RAG and GraphRAG under agentic search settings, where retrieval is performed iteratively during inference rather than appended as a one-shot input.